%% =====================================================================
%%  IEEE conference paper — Real-Time Explainable Hybrid CNN-LSTM IDS
%%
%%  Content source : docs/report/ieee_paper_content.md  (person-b-model)
%%  Every number below is taken from that file, which in turn cites
%%  committed results under model/results/.
%%
%%  Compile with pdfLaTeX (Overleaf default). No external files needed.
%%  Figures are deliberately excluded from this version.
%%
%%  WHAT YOU MUST FILL IN MANUALLY  (search for  %% TODO )
%%    1. Author e-mail addresses
%%    2. MU-IoT reference [8]
%%    3. Acknowledgment: add the MU-IoT credit
%% =====================================================================

\documentclass[conference]{IEEEtran}
\IEEEoverridecommandlockouts

\usepackage{cite}
\usepackage{amsmath,amssymb,amsfonts}
\usepackage{graphicx}
\usepackage{textcomp}
\usepackage{xcolor}
\usepackage{booktabs}
\usepackage{tabularx}
\usepackage{array}
\usepackage{url}
\usepackage{balance}
\usepackage{stfloats}

% nicer cell padding, no overfull boxes in narrow tables
\newcolumntype{Y}{>{\raggedright\arraybackslash}X}
\setlength{\tabcolsep}{4pt}
\renewcommand{\arraystretch}{1.08}

\begin{document}

\title{Real-Time Explainable Hybrid CNN-LSTM Intrusion Detection System with IoT-Based Data Collection and Cloud Deployment}

\author{%
\IEEEauthorblockN{Maanya B S, Ruthwik Sai Ganesh, Varshini D N, and Dr.~Madhu B R}
\IEEEauthorblockA{\textit{Department of Artificial Intelligence and Machine Learning}\\
\textit{Jyothy Institute of Technology}\\
Bengaluru, India\\
%% TODO: insert the four institutional e-mail addresses
[author e-mail addresses]}
}

\maketitle

%% ---------------------------------------------------------------------
\begin{abstract}
Deep-learning intrusion detection systems are usually reported as accuracy
figures on benchmark datasets, rarely as systems that act on live traffic.
This paper presents a hybrid CNN-LSTM intrusion detection and prevention
system that classifies sequences of network flows, explains its errors with
SHAP, and blocks attackers through the host firewall in real time. The model
is evaluated on CICIDS2017 with a chronological, leakage-controlled split and
on MU-IoT, real IoT traffic with four whole recordings held out. A random
forest trained on identical examples serves as a fair baseline: it is better
on CICIDS2017 (macro F1 0.780 vs.\ 0.586), whereas the CNN-LSTM is clearly
better on MU-IoT (held-out macro F1 0.676 vs.\ 0.404), where single flow
records are ambiguous and sequences resolve them. A per-class prevention
policy derives confidence thresholds from measured F1 and locks unreliable
classes to human review; on the CICIDS2017 test set it acts automatically on
96.6\% of attack windows and on 0.33\% of benign windows. A packet-to-flow
extractor reproduces CICFlowMeter-V3 features on 99.8\% of flows from the
original CICIDS2017 capture. In a live run on a shared Wi-Fi network, a real
port scan was detected about two seconds after its first probes and the
scanning device was blocked automatically, while another user's falsely
flagged device was refused. The limitations of a 2017-trained model on
present-day traffic are measured and reported.
\end{abstract}

\begin{IEEEkeywords}
Intrusion detection, intrusion prevention, CNN-LSTM, deep learning,
CICIDS2017, Internet of Things, explainable AI, SHAP, live packet capture.
\end{IEEEkeywords}

%% ---------------------------------------------------------------------
\section{Introduction}

Network intrusion detection systems (NIDS) based on machine learning are
commonly evaluated by accuracy on public datasets. Such figures can be
misleading: neighbouring flows in a capture are often near-identical, so
random train/test splits reward memorisation \cite{ref6}, \cite{ref7}, and
benign traffic dominates the data, so accuracy hides failures on rare
attacks. A detector that is to protect a network must also run on live
traffic and act on its decisions, which raises a question benchmark scores do
not answer: when may a model's prediction be trusted to block a device
automatically?

This work builds and evaluates a complete system around a hybrid CNN-LSTM
model that reads windows of consecutive flows. The contributions are:

\begin{enumerate}
\item \textbf{A leakage-controlled evaluation} of the CNN-LSTM on two
datasets, CICIDS2017 and MU-IoT, against a random-forest baseline trained and
tested on identical examples, showing where sequence modelling helps and
where it does not.
\item \textbf{An explainability analysis} that uses SHAP to separate three
causes of the model's errors on unseen IoT recordings.
\item \textbf{A prevention layer} in which each class's automatic-action
threshold follows its measured F1, unreliable classes are locked to human
review, and execution is bounded by a kill switch, an address allow-list and
one-rule-per-device bookkeeping.
\item \textbf{A live capture pipeline} whose packet-to-flow extractor is
validated against the original CICIDS2017 capture (99.8\% agreement), with a
live engine that detects and blocks in real time.
\item \textbf{A real-network demonstration} with its failure modes reported:
a port scan detected and blocked in about two seconds, and the false alarms a
2017-trained model raises on present-day IPv6 and QUIC traffic.
\end{enumerate}

\subsection*{Scope of the Built System}
The registered project title refers to IoT-based data collection and cloud
deployment. The scope of the present work is stated here explicitly to avoid
over-claiming. The IoT traffic used in this study is the external MU-IoT
dataset; no IoT sensor network was deployed or instrumented by the authors.
The system runs on the protected host itself, as a local backend and web
dashboard, and is not hosted in the cloud. All results reported in
Section~\ref{sec:results} were produced by that host-local deployment.

%% ---------------------------------------------------------------------
\section{Related Work}

\subsection{Deep Learning for Intrusion Detection}
Vinayakumar et al.~\cite{ref3} compared deep neural networks with classical
classifiers across several intrusion datasets and found deep models
competitive. Hybrid convolutional--recurrent models are a common choice for
flow sequences, convolution capturing local patterns between neighbouring
flows and recurrence capturing their development. This work uses such a
hybrid but measures it against a strong classical baseline on identical
examples.

\subsection{Datasets and Their Pitfalls}
CICIDS2017~\cite{ref1} provides labelled flows generated by
CICFlowMeter~\cite{ref2} from a testbed with benign users and fourteen attack
types. Engelen et al.~\cite{ref4} documented artefacts in CICFlowMeter's flow
construction, such as ``TCP appendix'' flows created after a connection's
FIN, and labelling errors. Sommer and Paxson~\cite{ref6} argued that network
intrusion detection differs from other machine-learning tasks in the cost of
errors and in the gap between training and deployment traffic; Arp et
al.~\cite{ref7} catalogued pitfalls such as data snooping and inappropriate
baselines. The evaluation design here follows these lessons: a chronological
split, a baseline on identical examples, and measurement on real traffic.

\subsection{Explainability}
SHAP~\cite{ref5} attributes a prediction to its input features with additive,
game-theoretic contributions. It is used here not to justify correct
predictions but to diagnose wrong ones.

%% ---------------------------------------------------------------------
\section{Datasets and Evaluation Design}

\subsection{CICIDS2017}
Fifteen classes (benign plus fourteen attacks) over five capture days.
Twenty features were selected by random-forest importance. Because
neighbouring flows are often near-identical, the data were split
chronologically within each (capture day, class) group: the first 80\% of
rows train, the remainder test. A pure day-level split was rejected because
it removed eight of fifteen classes from the test set. Windows of ten
consecutive flows never span two days, and the scaler is fitted on training
rows only. The test set holds 530{,}951 windows, 83\% benign.

\subsection{MU-IoT}
Real IoT traffic in seven classes (normal, DDoS, Injection, MiTM,
Password\_Hacking, Scan, Spyware), 3.94 million rows and 38
features~\cite{ref8}. Flow timestamps were removed because each recording
holds one class and its own time range, so timestamps would reveal the label.
Twenty-row windows never cross a recording block, a split or a gap in row
identifiers. Two test tracks are used: \emph{within} (later traffic of
recordings seen in training) and \emph{held-out} (four whole recordings never
seen: an hping3 UDP flood, a vulnerability scan, an MQTT dictionary attack
and an XSS attack). Headline scores count only test windows whose exact
sequence never appears in training.

\subsection{Baseline}
A random forest~\cite{ref10} (100 trees) is trained on single flows with the
same features and split, and scored on the last flow of each CNN-LSTM test
window, so both models are judged on identical examples and labels.

%% ---------------------------------------------------------------------
\section{Model}

The same architecture serves both datasets (Table~\ref{tab:arch}). Two
one-dimensional convolution blocks learn local patterns across neighbouring
flows; an LSTM~\cite{ref9} summarises their development across the window.

\begin{table*}[t]
\caption{CNN-LSTM Architecture}
\label{tab:arch}
\centering
\footnotesize
\begin{tabularx}{\textwidth}{@{}Ycc@{}}
\toprule
\textbf{Layer} & \textbf{CICIDS2017} & \textbf{MU-IoT}\\
\midrule
Input                               & 10 flows $\times$ 20 features & 20 flows $\times$ 38 features\\
Conv1D 64, kernel 3, ReLU; BatchNorm; MaxPool 2  & \checkmark & \checkmark\\
Conv1D 128, kernel 3, ReLU; BatchNorm; MaxPool 2 & \checkmark & \checkmark\\
LSTM, 64 units                      & \checkmark & \checkmark\\
Dense 64, ReLU                      & \checkmark & \checkmark\\
Dense, softmax                      & 15 classes & 7 classes\\
Parameters                          & 83{,}919   & 86{,}855\\
\bottomrule
\end{tabularx}
\end{table*}

Both models use Adam, sparse categorical cross-entropy, batch size 512 and
early stopping on validation loss (patience 5). For CICIDS2017, three
class-weighting schemes were compared on the same test set; square-root
balanced weights gave the best result (accuracy 0.984, macro F1 0.586,
against 0.883\,/\,0.428 for capped balanced weights). MU-IoT training draws a
fresh sample of up to 300{,}000 windows per class each epoch with
square-root-balanced weights capped at 10, in mixed precision.

%% ---------------------------------------------------------------------
\section{Prevention Layer}

\subsection{Policy}
Each class maps to an action, a severity and a confidence threshold.
Thresholds rise as a class's test F1 falls, from 0.55 (DDoS, DoS Hulk) to
0.78 (SSH-Patator). Below its threshold a prediction is held for human
review. Classes whose evidence is too weak never act automatically: on
CICIDS2017, Bot, Web Attack--Brute Force and XSS (F1 below 0.15) and SQL
Injection, Heartbleed and Infiltration (too few test examples); on MU-IoT,
MiTM and Spyware (no held-out recording). Measuring the policy changed it:
before the lock, all 17 automatic Bot actions on the CICIDS2017 test set hit
benign traffic.

\subsection{Execution}
An approved action becomes inbound and outbound Windows Firewall rules for
the attacker's address, added through \texttt{netsh advfirewall}. Execution
is bounded by independent checks: a global kill switch (off by default); an
allow-list of RFC~5737 documentation ranges plus single LAN devices listed
explicitly for a demonstration; a deny-list (loopback, broadcast,
unspecified); and an administrator check. One rule is kept per device
whatever class flags it, temporary actions expire on a timer, a startup sweep
recovers state after a crash, and a revoke operation lifts any block even
with the kill switch off. Rate limiting is approximated by a temporary
block, and input sanitisation for XSS stays a recommendation, since a network
firewall cannot perform it.

\subsection{Verification}
35 automated tests cover both policies and every refusal path; deliberately
breaking the engine in 26 ways made a test fail every time; an end-to-end
check on the real firewall passed 22 of 22 checks.

%% ---------------------------------------------------------------------
\section{Live Capture and Validation}

\subsection{Packet-to-Flow Extraction}
A new extractor groups packets from a live interface or a Wireshark capture
into bidirectional flows and computes the twenty features with
CICFlowMeter-V3's definitions, including the quirks the model learned: the
first packet counted twice in the packet-length statistics, sample standard
deviations, Ethernet padding counted as payload, whole-microsecond times, and
V3's flow termination rules. It reads IPv4 and IPv6 from raw bytes at about
68{,}000 packets per second.

\subsection{Validation Against the Original Data}
UNB's original Friday capture was passed through the extractor and every flow
compared with CICFlowMeter-V3's official row for it
(Table~\ref{tab:extractor}). The comparison found two differences, both
corrected: times must be whole microseconds, and a FIN that \emph{starts} a
flow (a TCP appendix~\cite{ref4}) does not end it.

\begin{table}[t]
\caption{Extractor Against CICFlowMeter-V3 (Original CICIDS2017 PCAP, 10 min)}
\label{tab:extractor}
\centering
\footnotesize
\begin{tabularx}{\columnwidth}{@{}Yr@{}}
\toprule
\textbf{Measure} & \textbf{Value}\\
\midrule
Flows matched one-to-one with official rows        & 5{,}957 of 7{,}332 (81\%)\\
All 20 features agree (at the CSV's printed precision) & 99.8\%\\
All 20 features agree (strictly, within $10^{-4}$) & 86.3\%\\
PortScan flows agreeing on all 20 features         & 49 of 49\\
Agreement if flows also end on RST                 & 98.9\%\\
\bottomrule
\end{tabularx}
\end{table}

\subsection{Live Engine}
Three threads separate capture, flow tracking and inference, so bursts of
packets never wait for the model. Once per second, if a flow has ended, the
engine classifies the ten most recently ended flows, selects the target (the
most frequent private IPv4 source that sent traffic to the host, never the
host itself or the gateway), and invokes the firewall engine itself, so
blocking does not depend on the dashboard. Live mode also ends flows on RST:
a scan flow (SYN followed by RST) has the same features either way, and a
scan is reported within seconds instead of after the 120\,s flow timeout. A
web dashboard shows predictions, executions and blocked devices with a revoke
control.

%% ---------------------------------------------------------------------
\section{Results}
\label{sec:results}

\subsection{Detection on CICIDS2017}

\begin{table}[t]
\caption{CICIDS2017, 530{,}951 Identical Test Examples}
\label{tab:cicids}
\centering
\footnotesize
\begin{tabularx}{\columnwidth}{@{}Yccc@{}}
\toprule
\textbf{Model} & \textbf{Accuracy} & \textbf{Weighted F1} & \textbf{Macro F1}\\
\midrule
CNN-LSTM (10-flow windows)   & 0.984 & 0.986 & 0.586\\
Random forest (single flows) & 0.991 & 0.991 & 0.780\\
\bottomrule
\end{tabularx}
\end{table}

Per class, the CNN-LSTM leads clearly only on DoS Hulk (F1 0.990 vs.\ 0.953).
Both models fail on Bot (0.056 and 0.031), which points to a feature or data
limit. The web attacks fail only for the CNN-LSTM (brute force 0.125 vs.\
0.688), so the signal exists in single flows and the windowed model loses it.
Large classes are detected well by both (DDoS 0.990, PortScan 0.971 for the
CNN-LSTM).

\subsection{Detection on MU-IoT}

\begin{table}[t]
\caption{MU-IoT, Unseen Windows Only}
\label{tab:muiot}
\centering
\footnotesize
\begin{tabularx}{\columnwidth}{@{}Ycc@{}}
\toprule
\textbf{Measure} & \textbf{CNN-LSTM} & \textbf{Random forest}\\
\midrule
Within: accuracy / macro F1             & 0.964 / 0.936 & 0.904 / 0.816\\
Held-out: accuracy / macro F1           & 0.680 / 0.676 & 0.407 / 0.404\\
Normal traffic flagged as attack (within) & 9.3\% & 17.2\%\\
Held-out attacks called normal          & 0.09\% & 2.5\%\\
\bottomrule
\end{tabularx}
\end{table}

The CNN-LSTM is clearly better, most of all on recordings never seen. It
almost never misses an attack; its held-out errors concern the attack's name,
as unfamiliar tools are mapped to the nearest familiar class. The reason
sequences help here is measurable: with timestamps removed, 71.8\% of single
Scan test records also appear in training under another label, against 0.0\%
of twenty-row Scan windows.

\subsection{Error Analysis with SHAP}
Expected-gradient SHAP values (validated against
\texttt{shap.GradientExplainer}, $r = 0.990$) were averaged per error group
and compared, by cosine similarity, with correctly classified groups
(Table~\ref{tab:shap}).

\begin{table*}[t]
\caption{Causes of the Main MU-IoT Errors}
\label{tab:shap}
\centering
\footnotesize
\begin{tabularx}{\textwidth}{@{}YccY@{}}
\toprule
\textbf{Error} & \textbf{Similarity to predicted class} &
\textbf{Similarity to true class} & \textbf{Cause}\\
\midrule
normal $\rightarrow$ Spyware & 0.87 & 0.29 & Systematic resemblance\\
hping3 flood $\rightarrow$ Scan & 0.59 & $-0.64$ & Real resemblance to scans\\
vulnerability scan $\rightarrow$ Password\_Hacking & 0.20 & $-0.05$ &
Extrapolation beyond training range\\
\bottomrule
\end{tabularx}
\end{table*}

The first error cannot be fixed by a threshold, which is why Spyware is locked
to human review; the second reflects a gap in training coverage; the third is
driven by one feature whose values lie far outside the training range.

\subsection{Prevention Policy}

\begin{table*}[t]
\caption{Policy Outcomes on All Test Windows}
\label{tab:policy}
\centering
\footnotesize
\begin{tabularx}{\textwidth}{@{}Yccc@{}}
\toprule
\textbf{Dataset} & \textbf{Benign windows auto-actioned} &
\textbf{Attacks handled automatically} & \textbf{Attacks missed}\\
\midrule
CICIDS2017                    & 0.33\% & 96.6\% & 2.4\%\\
MU-IoT, within recordings     & 1.6\%  & 94.1\% & 1.5\%\\
MU-IoT, held-out recordings   & ---    & 94.4\% & 0.08\%\\
\bottomrule
\end{tabularx}
\end{table*}

\subsection{Real Traffic and the Live Run}
Three captures were recorded on a phone hotspot with a laptop as the protected
host and a tablet as the second device. In a port scan (\texttt{nmap -sT},
1{,}051 probes in 7\,s), 98\% of the windows pointing at the tablet were
classified PortScan. In laptop browsing, 11.1\% of windows were not benign and
7.6\% recommended an automatic action, none of them with a target that could
be blocked; a second, ten-minute browsing capture reached 29.9\% non-benign
windows. While the tablet repeatedly loaded a page on the laptop, its windows
were read as Bot and held for review, never executed.

The live run took place on a home Wi-Fi network used daily by 20--30 people,
with the owner's permission and the tablet as the only listed device:

\begin{itemize}
\item A real port scan was classified PortScan (confidence 0.99) about two
seconds after its first probes; the tablet was blocked with one rule, and
twelve further detections were answered ``already blocked''.
\item The laptop's services became unreachable from the tablet; revoking the
rule restored access. Three full scan, block and revoke cycles were run.
\item Another user's device was classified as a slow-HTTP DoS (confidence
0.996) from its SSDP discovery broadcasts; the engine refused all ten attempts
because the device was not listed. Since a broadcasting device never addresses
the host, targets are now required to have sent traffic to it.
\end{itemize}

Single-window inference takes 1.48\,ms on a Ryzen~3 laptop CPU, and batched
inference about 70{,}700 windows per second.

%% ---------------------------------------------------------------------
\section{Discussion and Limitations}

\begin{itemize}
\item \textbf{Sequences are not a universal improvement.} They help where
single flows are ambiguous (MU-IoT) and lose to a single-flow forest where
flows are distinctive (CICIDS2017).
\item \textbf{Dataset shift.} The 2017-trained model raises
recommendation-only false alarms on present-day IPv6 and QUIC traffic
(11--30\% of browsing windows), reads repeated polling as Bot, and recognises
a scan only when the victim answers it (Windows' stealth mode had to be
disabled for the demonstration). The safety layers kept all of these from
being executed.
\item \textbf{Scope of the live validation.} One live attack type (a TCP
connect scan) was run, on one network, from one device. Floods, brute-force
and web attacks require a victim, vulnerable services or several attacking
machines on an isolated laboratory network, which were not available; these
classes were evaluated on recorded data only. Blocking follows detection by
about two seconds, so a fast scan has already probed most ports by then.
\item \textbf{Evaluation caveats.} Feature selection used a sample that
includes some later test rows; thresholds come from test-set F1 without a
separate calibration set; each model was trained once, so scores carry no
error margin; the extractor was validated on one ten-minute slice of one day.
\item \textbf{Prevention is deliberately limited} (IPS-lite): it blocks only
listed devices or documentation addresses, behind a kill switch.
\end{itemize}

%% ---------------------------------------------------------------------
\section{Conclusion and Future Work}

A hybrid CNN-LSTM intrusion detection and prevention system was evaluated
with leakage-controlled splits and a baseline on identical examples, and
taken from benchmark data to live traffic. Its strengths are specific: it is
clearly better than a random forest on ambiguous IoT traffic and on unseen
attack recordings, but not on CICIDS2017. Deriving each class's
automatic-action threshold from its measured F1, and locking unreliable
classes to review, kept automatic action on benign traffic at 0.33\% on the
CICIDS2017 test set, and the live run showed real attacks blocked within
seconds while false alarms on other devices were refused. Future work
includes fine-tuning on present-day benign traffic (IPv6, QUIC), evaluating
further attack types live on an isolated laboratory network, a separate
calibration set for the thresholds, and repeated training runs to report
score variance.

%% ---------------------------------------------------------------------
\section*{Acknowledgment}

The authors thank Dr.~Madhu B R for guidance and the Canadian Institute for
Cybersecurity (University of New Brunswick) for the CICIDS2017 dataset.
%% TODO: add the MU-IoT dataset acknowledgment here.

%% ---------------------------------------------------------------------
\balance
\begin{thebibliography}{10}

\bibitem{ref1} I.~Sharafaldin, A.~H.~Lashkari, and A.~A.~Ghorbani, ``Toward
generating a new intrusion detection dataset and intrusion traffic
characterization,'' in \emph{Proc. 4th Int. Conf. Information Systems Security
and Privacy (ICISSP)}, 2018, pp. 108--116.

\bibitem{ref2} A.~H.~Lashkari, G.~Draper-Gil, M.~S.~I.~Mamun, and
A.~A.~Ghorbani, ``Characterization of Tor traffic using time based
features,'' in \emph{Proc. 3rd Int. Conf. Information Systems Security and
Privacy (ICISSP)}, 2017, pp. 253--262.

\bibitem{ref3} R.~Vinayakumar, M.~Alazab, K.~P.~Soman, P.~Poornachandran,
A.~Al-Nemrat, and S.~Venkatraman, ``Deep learning approach for intelligent
intrusion detection system,'' \emph{IEEE Access}, vol.~7, pp. 41525--41550,
2019.

\bibitem{ref4} G.~Engelen, V.~Rimmer, and W.~Joosen, ``Troubleshooting an
intrusion detection dataset: The CICIDS2017 case study,'' in \emph{Proc. IEEE
Security and Privacy Workshops (SPW)}, 2021, pp. 7--12.

\bibitem{ref5} S.~M.~Lundberg and S.-I.~Lee, ``A unified approach to
interpreting model predictions,'' in \emph{Advances in Neural Information
Processing Systems 30 (NeurIPS)}, 2017.

\bibitem{ref6} R.~Sommer and V.~Paxson, ``Outside the closed world: On using
machine learning for network intrusion detection,'' in \emph{Proc. IEEE
Symposium on Security and Privacy}, 2010, pp. 305--316.

\bibitem{ref7} D.~Arp \emph{et al.}, ``Dos and don'ts of machine learning in
computer security,'' in \emph{Proc. 31st USENIX Security Symposium}, 2022.

%% TODO: replace the entry below with the real MU-IoT dataset reference.
\bibitem{ref8} MU-IoT dataset, Manipur University, IoT intrusion traffic
recordings. [Reference to be completed by the data team.]

\bibitem{ref9} S.~Hochreiter and J.~Schmidhuber, ``Long short-term memory,''
\emph{Neural Computation}, vol.~9, no.~8, pp. 1735--1780, 1997.

\bibitem{ref10} L.~Breiman, ``Random forests,'' \emph{Machine Learning},
vol.~45, no.~1, pp. 5--32, 2001.

\end{thebibliography}

\end{document}
